% ----------------------------------------------------------
% Justificativa
% ----------------------------------------------------------
\chapter{Justificativa} \label{cha:justificativa}


O mercado brasileiro de jogos digitais tem se tornado um ecossistema cada vez mais diversificado, acompanhando o desenvolvimento da indústria de jogos no país \cite{abragames2023}. Nesse contexto, a comunicação entre desenvolvedores e jogadores assume um papel relevante, pois permite a coleta de \textit{feedbacks} ao longo do desenvolvimento e a realização de testes antecipados, contribuindo para a identificação de problemas e o aprimoramento dos projetos \cite{bain2025}.

Entretanto, a multiplicidade de canais de comunicação aumenta a complexidade do processo de triagem e da análise dessas contribuições, além de favorecer a perda de contexto decorrente da descentralização e do fluxo desordenado de informações \cite{xiao2025}. Dessa forma, a ausência de um ambiente unificado pode dificultar a localização, a organização e o acompanhamento dos \textit{feedbacks} dos jogadores, tornando menos eficiente a identificação de informações relevantes em meio ao volume de dados recebidos.

Nesse cenário, uma interface que promova a centralização e a organização desses conteúdos pode facilitar o acesso à comunidade e a interação entre os indivíduos. Para isso, aspectos como a usabilidade, a experiência de usuário (\textit{User Experience} -- UX), a acessibilidade e a comunicabilidade são essenciais para proporcionar uma interação satisfatória com a plataforma, permitindo que os usuários naveguem entre as informações e realizem suas tarefas de maneira adequada \cite{barbosa2010}.

Além disso, para validar escolhas de \textit{design} e assegurar que a plataforma atenda às necessidades do público, a realização de testes com usuários se torna uma etapa indispensável. Segundo \citeonline{nielsen1993}, o teste de usabilidade com usuários reais é o método mais fundamental por fornecer informações diretas sobre o comportamento e as dificuldades de interação. 

Portanto, a relevância deste trabalho está na proposta de centralização de \textit{feedbacks} relacionados a jogos digitais por meio de uma interface que facilite a organização e a consulta dessas contribuições. No âmbito acadêmico, o estudo pode contribuir para discussões sobre o desenvolvimento e a avaliação de interfaces voltadas à participação de comunidades digitais, enquanto, no âmbito social, busca favorecer a comunicação entre jogadores e desenvolvedores, oferecendo um ambiente que facilite essa interação.